%!TEX root = ../report.tex
\documentclass[report.tex]{subfiles}
\begin{document}
    \chapter{Results and Discussion} \label{chap:resultsanddiscussion}

    The evaluation of the proposed method using the experiments and evaluation metrics discussed in the previous chapter is presented in this chapter. The performance comparison of the proposed \gls{cvo} approach and the baseline odometry against the Vicon-based ground truth is discussed. This chapter is organized based on the terrain types used in the experiments. The final part of this chapter presents the discussion of the results obtained from the experiments. 
    
    \section{Flat Floor Terrain Experiments} \label{sec:flatFloorTerrainResults}   
    
    The evaluation results from the datasets presented in Section~\ref{subsec:exp_flatFloorTerrain} are discussed in this section. A total of 53 datasets were collected on flat floor terrain experiments. Since presenting all datasets would make the chapter repetitive, a set of representative datasets is selected for detailed discussion. The combined plots show the ground truth, baseline odometry, and the contact velocity odometry data together in a single plot. The following subsections present selected combined 2D trajectory plots, \texttt{z}-position comparison plots, and translational and rotational \gls{rpe} plots for the baseline odometry and contact velocity odometry. Where relevant, translational and rotational \gls{ape} plots are also presented. For each selected dataset, a summary table is provided with the final pose error, translational and rotational drift rates, and travel distance error.
    
    The representative datasets are selected to include different trajectory types, such as straight-line motion, circular motion, and complex trajectories. This selection helps to discuss the performance of the proposed method on simple and complex trajectories.  Unless stated otherwise, the same color coding is used across all plots. The ground truth trajectory is shown in blue, the baseline odometry trajectory is shown in green, and the contact velocity odometry trajectory is shown in red. The yaw orientations are shown as heading arrows along the trajectories using the same respective colors.
    
    \subsection{Straight-Line Motion}
    
    The straight-line motion datasets include the experiments where the rover moves in a straight-line without any turns. Three straight-line motion types are included on flat floor terrain experiments. The straight-line motion datasets are: \texttt{DATA\_02\_YY}, \texttt{DATA\_03\_YY}, and \texttt{DATA\_04\_YY}, where \texttt{YY} represents the iteration number. The evaluation results of these datasets are presented in the following subsections.
    
    \subsubsection{Straight-Line Motion Dataset 1: \texttt{DATA\_02\_03}} \label{subsubsec:straightLineD1_02_03}
    
    In this dataset, the Coyote~3 rover moves forward in a straight-line for a distance of approximately $5~\mathrm{m}$. The trajectory of the rover does not include any turns or complex maneuvers. The maximum commanded velocity of the rover is $0.2~\mathrm{m/s}$. 
    The X--Y trajectory comparison is shown in Figure~\ref{fig:straight_line_motion_dataset_02_03_trajectory}. The \texttt{z}-position comparison plot is shown in Figure~\ref{fig:straight_line_motion_dataset_02_03_z_position}. The color coding is the same as in the X--Y trajectory plot. The translational and rotational \gls{rpe} plots are shown in Figures~\ref{fig:straight_line_motion_dataset_02_03_RPE_translational} and~\ref{fig:straight_line_motion_dataset_02_03_RPE_rotational}, respectively. Metrics such as travel distance error, final position error, translational and rotational drift rates are summarized in Table~\ref{tab:odometry_metrics_data-02-03}.
            
    \begin{figure}[!htbp]
        \centering
        \includegraphics[width=0.75\textwidth]{media/results/01_flat_floor_terrain/DATA_02_03/Trajectory_With_Heading_vicon_ground_truth_vs_contact_velocity_odometry_vs_baseline_odometry_DATA_02_03.png}
        \caption{X--Y trajectory plot of the straight-line motion dataset \texttt{DATA\_02\_03} on the flat floor terrain. The yaw orientations are shown as heading arrows along the trajectories. The plot compares the ground truth, baseline odometry, and contact velocity odometry.}
        \label{fig:straight_line_motion_dataset_02_03_trajectory}
    \end{figure}
    
    \begin{figure}[!htbp]
        \centering
        \includegraphics[width=\textwidth]{media/results/01_flat_floor_terrain/DATA_02_03/z_position_comparison_vicon_ground_truth_vs_contact_velocity_odometry_vs_baseline_odometry_DATA_02_03.png}
        \caption{\texttt{z}-position plot of the straight-line motion dataset \texttt{DATA\_02\_03} on flat floor terrain. The plot compares the ground truth, baseline odometry, and contact velocity odometry along the \texttt{z} axis.}
        \label{fig:straight_line_motion_dataset_02_03_z_position}
    \end{figure}
    
    \begin{figure}[!htbp]
        \centering
        \includegraphics[width=\textwidth]{media/results/01_flat_floor_terrain/DATA_02_03/RPE_translational_comparison_contact_velocity_odometry_vs_baseline_odometry_DATA_02_03.png}
        \caption{Translational \gls{rpe} plot of the straight-line motion dataset \texttt{DATA\_02\_03} on the flat floor terrain. The plot compares the translational \gls{rpe} of the baseline odometry and contact velocity odometry against the ground truth.}
        \label{fig:straight_line_motion_dataset_02_03_RPE_translational}
    \end{figure}
    
    \begin{figure}[!htbp]
        \centering
        \includegraphics[width=\textwidth]{media/results/01_flat_floor_terrain/DATA_02_03/RPE_rotational_comparison_contact_velocity_odometry_vs_baseline_odometry_DATA_02_03.png}
        \caption{Rotational \gls{rpe} plot of the straight-line motion dataset \texttt{DATA\_02\_03} on the flat floor terrain. The plot compares the rotational \gls{rpe} of the baseline odometry and contact velocity odometry against the ground truth.}
        \label{fig:straight_line_motion_dataset_02_03_RPE_rotational}
    \end{figure}

    
    \begin{table}[ht]
    \centering
    \caption{Odometry evaluation metrics for DATA\_02\_03.}
    \label{tab:odometry_metrics_data-02-03}
        \begin{tabular}{lll}
        \hline
        Metric & Baseline odometry & Contact velocity odometry \\
        \hline
        Travel distance error & -0.18~m & -0.76~m \\
        Final position error & 0.18~m & 0.89~m \\
        RPE translational drift RMSE & 0.05~m/m & 0.14~m/m \\
        RPE rotational drift RMSE & 1.04~$^\circ$/m & 2.69~$^\circ$/m \\
        \hline
        \end{tabular}
    \end{table}
    
    \todo[inline]{include additional metrics to the starting straight line dataset tables, CHECK THEM ALL}
    
    From Figure~\ref{fig:straight_line_motion_dataset_02_03_trajectory}, it can be observed that the baseline odometry trajectory closely follows the ground truth, while the \gls{cvo} trajectory shows an initial yaw deviation, which is reflected in the deviation of the X--Y trajectory.
    
    From Table~\ref{tab:odometry_metrics_data-02-03}, it can be observed that the \gls{cvo} approach has a higher travel distance error, final position error, and translational and rotational drift rates compared to the baseline odometry. A negative travel distance error indicates that the estimated travel distance is smaller than the ground truth travel distance. The \gls{rmse} of the translational and rotational \gls{rpe} drift is higher for the \gls{cvo} compared to the baseline odometry, as shown in Figures~\ref{fig:straight_line_motion_dataset_02_03_RPE_translational} and~\ref{fig:straight_line_motion_dataset_02_03_RPE_rotational}. Notably, the \texttt{z}-position plot shown in Figure~\ref{fig:straight_line_motion_dataset_02_03_z_position} indicates that the \gls{cvo} approach tracks the ground truth better than the baseline odometry. Based on the evaluation results, when the rover moves in a straight-line forward without any turns or complex maneuvers, the baseline odometry performs better than the \gls{cvo} approach. 
    
    \subsubsection{Straight-Line Motion Dataset 2: \texttt{DATA\_03\_01}} \label{subsubsec:straightLineDD2_03_01}
    
    In this dataset, the Coyote~3 rover moves backward in a straight-line for a distance of approximately $5~\mathrm{m}$. The trajectory of the rover does not include any turns or complex maneuvers. The maximum commanded velocity of the rover is $0.2~\mathrm{m/s}$. The X--Y trajectory comparison and the \texttt{z}-position comparison for \texttt{DATA\_03\_01} are shown in Figures~\ref{fig:straight_line_motion_dataset_03_01} and~\ref{fig:straight_line_motion_dataset_03_01_z_position}, respectively. The translational and rotational \gls{rpe} plots are shown in Figures~\ref{fig:straight_line_motion_dataset_03_01_RPE_translational} and~\ref{fig:straight_line_motion_dataset_03_01_RPE_rotational}, respectively. Some key evaluation metrics are summarized in Table~\ref{tab:odometry_metrics_data-03-01}.
    
    \begin{figure}[!htbp]
        \centering
        \includegraphics[width=0.75\textwidth]{media/results/01_flat_floor_terrain/DATA_03_01/Trajectory_With_Heading_vicon_ground_truth_vs_contact_velocity_odometry_vs_baseline_odometry_DATA_03_01.png}
        \caption{X--Y trajectory plot of the straight-line motion dataset \texttt{DATA\_03\_01} on the flat floor terrain. The yaw orientations are shown as heading arrows along the trajectories. The plot compares the ground truth, baseline odometry, and contact velocity odometry.}
        \label{fig:straight_line_motion_dataset_03_01}
    \end{figure}
    
    \begin{figure}[!htbp]
        \centering
        \includegraphics[width=\textwidth]{media/results/01_flat_floor_terrain/DATA_03_01/z_position_comparison_vicon_ground_truth_vs_contact_velocity_odometry_vs_baseline_odometry_DATA_03_01.png}
        \caption{\texttt{z}-position plot of the straight-line motion dataset \texttt{DATA\_03\_01} on flat floor terrain. The plot compares the ground truth, baseline odometry, and contact velocity odometry along the \texttt{z} axis.}
        \label{fig:straight_line_motion_dataset_03_01_z_position}
    \end{figure}
    
    \begin{figure}[!htbp]
        \centering
        \includegraphics[width=\textwidth]{media/results/01_flat_floor_terrain/DATA_03_01/RPE_translational_comparison_contact_velocity_odometry_vs_baseline_odometry_DATA_03_01.png}
        \caption{Translational \gls{rpe} plot of the straight-line motion dataset \texttt{DATA\_03\_01} on the flat floor terrain. The plot compares the translational \gls{rpe} of the baseline odometry and contact velocity odometry against the ground truth.}
        \label{fig:straight_line_motion_dataset_03_01_RPE_translational}
    \end{figure}
    
    \begin{figure}[!htbp]
        \centering
        \includegraphics[width=\textwidth]{media/results/01_flat_floor_terrain/DATA_03_01/RPE_rotational_comparison_contact_velocity_odometry_vs_baseline_odometry_DATA_03_01.png}
        \caption{Rotational \gls{rpe} plot of the straight-line motion dataset \texttt{DATA\_03\_01} on the flat floor terrain. The plot compares the rotational \gls{rpe} of the baseline odometry and contact velocity odometry against the ground truth.}
        \label{fig:straight_line_motion_dataset_03_01_RPE_rotational}
    \end{figure}

    \begin{table}[ht]
    \centering
    \caption{Odometry evaluation metrics for DATA\_03\_01.}
    \label{tab:odometry_metrics_data-03-01}
    \begin{tabular}{lll}
    \hline
    Metric & Baseline odometry & Contact velocity odometry \\
    \hline
    Travel distance error & -0.12~m & -0.68~m \\
    Final position error & 0.46~m & 0.83~m \\
    RPE translational drift RMSE & 0.07~m/m & 0.14~m/m \\
    RPE rotational drift RMSE & 3.42~$^\circ$/m & 3.79~$^\circ$/m \\
    \hline
    \end{tabular}
    \end{table}
    
    \begin{figure}[!htbp]
        \centering
        \includegraphics[width=\textwidth]{media/results/01_flat_floor_terrain/DATA_03_01/yaw_orientation_comparison_vicon_ground_truth_vs_contact_velocity_odometry_vs_baseline_odometry_DATA_03_01.png}
        \caption{Yaw orientation comparison for the straight-line motion dataset \texttt{DATA\_03\_01} on the flat floor terrain. The plot compares the yaw orientations of the baseline odometry and contact velocity odometry against the ground truth.}
        \label{fig:straight_line_motion_dataset_03_01_yaw_orientation}
    \end{figure}

    From Figure~\ref{fig:straight_line_motion_dataset_03_01}, it can be seen that both the baseline odometry and \gls{cvo} trajectories deviate from the ground truth. Similar to the forward motion, the travel distance estimation of the baseline odometry is closer to the ground truth.
    
    The \gls{cvo} underestimates the travel distance, which can be seen in the X--Y trajectory plot and in the \texttt{Travel distance error} row in Table~\ref{tab:odometry_metrics_data-03-01}. The \gls{rmse} of translational drift is consistent with the forward motion \texttt{DATA\_02\_03} dataset, as shown in Figure~\ref{fig:straight_line_motion_dataset_03_01_RPE_translational} and Table~\ref{tab:odometry_metrics_data-03-01}.
    
    From Figure~\ref{fig:straight_line_motion_dataset_03_01_yaw_orientation}, the yaw error at the start of the trajectory causes both the baseline odometry and \gls{cvo} trajectories to deviate from the ground truth. This yaw estimation becomes better once the rover starts moving around $18~\mathrm{s}$, this can be seen in both Figure~\ref{fig:straight_line_motion_dataset_03_01_yaw_orientation} and Figure~\ref{fig:straight_line_motion_dataset_03_01_RPE_rotational}. Similar to the forward motion, when the trajectory is simple and does not include any turns or complex maneuvers, the baseline odometry performs better than the \gls{cvo} approach.

    \subsubsection{Straight-Line Motion Dataset 3: \texttt{DATA\_04\_02}}  \label{subsubsec:straightLineD3_04_02}
    
    In this dataset, the Coyote~3 rover moves forward in a straight-line for approximately $5~\mathrm{m}$, stops for 4--5 seconds, and then moves backward for approximately $5~\mathrm{m}$. The maximum commanded velocity of the rover is $0.2~\mathrm{m/s}$ when moving. This motion is grouped under straight-line motion since the rover moves in a straight-line without any turns or complex maneuvers. 
    
    The X--Y trajectory comparison and the \texttt{z}-position comparison for \texttt{DATA\_04\_02} are shown in Figures~\ref{fig:straight_line_motion_dataset_04_02_trajectory} and~\ref{fig:straight_line_motion_dataset_04_02_z_position}, respectively. The translational and rotational \gls{rpe} plots are shown in Figures~\ref{fig:straight_line_motion_dataset_04_02_RPE_translational} and~\ref{fig:straight_line_motion_dataset_04_02_RPE_rotational}, respectively. Table~\ref{tab:odometry_metrics_data-04-02} summarizes some key evaluation metrics for this dataset.
    

    \begin{figure}[!htbp]
        \centering
        \includegraphics[width=0.75\textwidth]{media/results/01_flat_floor_terrain/DATA_04_02/Trajectory_With_Heading_vicon_ground_truth_vs_contact_velocity_odometry_vs_baseline_odometry_DATA_04_02.png}
        \caption{X--Y trajectory comparison for the straight-line motion dataset \texttt{DATA\_04\_02} on the flat floor terrain. The plot compares the trajectories of the baseline odometry and contact velocity odometry against the ground truth.}
        \label{fig:straight_line_motion_dataset_04_02_trajectory}
    \end{figure}
    
    \begin{figure}[!htbp]
        \centering
        \includegraphics[width=\textwidth]{media/results/01_flat_floor_terrain/DATA_04_02/z_position_comparison_vicon_ground_truth_vs_contact_velocity_odometry_vs_baseline_odometry_DATA_04_02.png}
        \caption{\texttt{z}-position plot of the straight-line motion dataset \texttt{DATA\_04\_02} on flat floor terrain. The plot compares the ground truth, baseline odometry, and contact velocity odometry along the \texttt{z} axis.}
        \label{fig:straight_line_motion_dataset_04_02_z_position}
    \end{figure}
    
    \begin{figure}[!htbp]
        \centering
        \includegraphics[width=\textwidth]{media/results/01_flat_floor_terrain/DATA_04_02/RPE_translational_comparison_contact_velocity_odometry_vs_baseline_odometry_DATA_04_02.png}
        \caption{Translational \gls{rpe} comparison for the straight-line motion dataset \texttt{DATA\_04\_02} on the flat floor terrain. The plot compares the translational \gls{rpe} of the baseline odometry and contact velocity odometry against the ground truth.}
        \label{fig:straight_line_motion_dataset_04_02_RPE_translational}
    \end{figure}
    
    \begin{figure}[!htbp]
        \centering
        \includegraphics[width=\textwidth]{media/results/01_flat_floor_terrain/DATA_04_02/RPE_rotational_comparison_contact_velocity_odometry_vs_baseline_odometry_DATA_04_02.png}
        \caption{Rotational \gls{rpe} comparison for the straight-line motion dataset \texttt{DATA\_04\_02} on the flat floor terrain. The plot compares the rotational \gls{rpe} of the baseline odometry and contact velocity odometry against the ground truth.}
        \label{fig:straight_line_motion_dataset_04_02_RPE_rotational}
    \end{figure}      
        
    \begin{table}[ht]
    \centering
    \caption{Odometry evaluation metrics for DATA\_04\_02.}
    \label{tab:odometry_metrics_data-04-02}
    \begin{tabular}{lll}
    \hline
    Metric & Baseline odometry & Contact velocity odometry \\
    \hline
    Travel distance error & -0.54~m & -1.40~m \\
    Final position error & 0.41~m & 0.45~m \\
    RPE translational drift RMSE & 0.07~m/m & 0.16~m/m \\
    RPE rotational drift RMSE & 2.11~$^\circ$/m & 2.60~$^\circ$/m \\
    \hline
    \end{tabular}
    \end{table}

    This dataset combines the forward and backward straight-line motion that is discussed in the previous two datasets. Similar to \texttt{DATA\_03\_01} backward straight-line dataset, the \gls{cvo} X--Y trajectory deviates from the ground truth and closely matches the baseline odometry trajectory. This is shown in the X--Y trajectory plot in Figure~\ref{fig:straight_line_motion_dataset_04_02_trajectory}. 
    
    The travel distance error of \gls{cvo} is higher than the baseline odometry, as shown in Table~\ref{tab:odometry_metrics_data-04-02}. This is the case for all the straight-line trajectories on flat floor terrain. Notably, the final position error of \gls{cvo} is closer to the baseline odometry final position error than in the previous two datasets.
    
    \subsection{Circular Motion}
    
    The circular motion datasets include the experiments where the rover moves in a circular trajectory. The circular motion datasets are: \texttt{DATA\_09\_YY} and \texttt{DATA\_10\_YY}, where \texttt{YY} represents the iteration number. The evaluation results of these datasets are presented in the following subsections. 
    
    \subsubsection{Circular Motion Dataset 1: \texttt{DATA\_09\_03}} \label{subsubsec:circularMotionD1_09_03}
    
    In this dataset, the Coyote~3 rover moves in a circular trajectory with a diameter of approximately $2.5-2.8~\mathrm{m}$. The circular motion is performed in a counter-clockwise direction. This rotation direction can be seen in the X--Y trajectory plot in Figure~\ref{fig:circular_motion_dataset_09_03_trajectory}. The trajectory of the rover does not include any complex maneuvers.
    
    The maximum commanded velocity of the rover is $0.2~\mathrm{m/s}$. The X--Y trajectory comparison and the \texttt{z}-position comparison for \texttt{DATA\_09\_03} are shown in Figures~\ref{fig:circular_motion_dataset_09_03_trajectory} and~\ref{fig:circular_motion_dataset_09_03_z_position}, respectively. The translational and rotational \gls{rpe} plots are shown in Figures~\ref{fig:circular_motion_dataset_09_03_RPE_translational} and~\ref{fig:circular_motion_dataset_09_03_RPE_rotational}, respectively. Table~\ref{tab:odometry_metrics_data-09-03} summarizes some key evaluation metrics for this dataset.

    \begin{figure}[!htbp]
        \centering
        \includegraphics[width=0.75\textwidth]{media/results/01_flat_floor_terrain/DATA_09_03/Trajectory_With_Heading_vicon_ground_truth_vs_contact_velocity_odometry_vs_baseline_odometry_DATA_09_03.png}
        \caption{X--Y trajectory comparison for the circular motion dataset \texttt{DATA\_09\_03} on the flat floor terrain. The plot compares the trajectories of the baseline odometry and contact velocity odometry against the ground truth.}
        \label{fig:circular_motion_dataset_09_03_trajectory}
    \end{figure}
    
    \begin{figure}[!htbp]
        \centering
        \includegraphics[width=\textwidth]{media/results/01_flat_floor_terrain/DATA_09_03/z_position_comparison_vicon_ground_truth_vs_contact_velocity_odometry_vs_baseline_odometry_DATA_09_03.png}
        \caption{\texttt{z}-position plot of the circular motion dataset \texttt{DATA\_09\_03} on flat floor terrain. The plot compares the ground truth, baseline odometry, and contact velocity odometry along the \texttt{z} axis.}
        \label{fig:circular_motion_dataset_09_03_z_position}
    \end{figure}
    
    \begin{figure}[!htbp]
        \centering
        \includegraphics[width=\textwidth]{media/results/01_flat_floor_terrain/DATA_09_03/RPE_translational_comparison_contact_velocity_odometry_vs_baseline_odometry_DATA_09_03.png}
        \caption{Translational \gls{rpe} comparison for the circular motion dataset \texttt{DATA\_09\_03} on the flat floor terrain. The plot compares the translational \gls{rpe} of the baseline odometry and contact velocity odometry against the ground truth.}
        \label{fig:circular_motion_dataset_09_03_RPE_translational}
    \end{figure}
    
    \begin{figure}[!htbp]
        \centering
        \includegraphics[width=\textwidth]{media/results/01_flat_floor_terrain/DATA_09_03/RPE_rotational_comparison_contact_velocity_odometry_vs_baseline_odometry_DATA_09_03.png}
        \caption{Rotational \gls{rpe} comparison for the circular motion dataset \texttt{DATA\_09\_03} on the flat floor terrain. The plot compares the rotational \gls{rpe} of the baseline odometry and contact velocity odometry against the ground truth.}
        \label{fig:circular_motion_dataset_09_03_RPE_rotational}
    \end{figure}      
    
    \begin{table}[ht]
    \centering
    \caption{Odometry evaluation metrics for DATA\_09\_03.}
    \label{tab:odometry_metrics_data-09-03}
    \begin{tabular}{lll}
    \hline
    Metric & Baseline odometry & Contact velocity odometry \\
    \hline
    Travel distance error & -0.06~m & -1.06~m \\
    Final position error & 0.15~m & 0.16~m \\
    Yaw RMSE & 4.89~$^\circ$ & 7.71~$^\circ$ \\
    Z position RMSE & 0.01~m & 0.04~m \\
    RPE translational drift RMSE & 0.16~m/m & 0.16~m/m \\
    RPE rotational drift RMSE & 5.17~$^\circ$/m & 4.92~$^\circ$/m \\
    \hline
    \end{tabular}
    \end{table}
    
    The motion in this dataset is purely circular and does not include straight-line motion. From the X--Y trajectory plot in Figure~\ref{fig:circular_motion_dataset_09_03_trajectory}, it can be observed that both the baseline odometry and \gls{cvo} trajectories follow the shape of the ground truth trajectory. The translational \gls{rpe} drift \gls{rmse} is the same for both the baseline odometry and \gls{cvo}. The rotational \gls{rpe} drift \gls{rmse} is slightly better for the \gls{cvo} compared to the baseline odometry. This can be inferred from Table~\ref{tab:odometry_metrics_data-09-03} and Figures~\ref{fig:circular_motion_dataset_09_03_RPE_translational} and~\ref{fig:circular_motion_dataset_09_03_RPE_rotational}. The \texttt{z}-position plot in Figure~\ref{fig:circular_motion_dataset_09_03_z_position} shows that the \gls{cvo} follows the variations in the ground truth better than the baseline odometry. Over time the \gls{cvo} \texttt{z}-position estimation drifts away from the ground truth. This increases the \texttt{z}-position \gls{rmse}. Unlike the straight-line motion datasets, the final position error of the \gls{cvo} is closer to the baseline odometry final position error. The travel distance errors are both negative, meaning both the \gls{cvo} and the baseline odometry underestimated the true travel distance. Compared to the straight-line datasets, the baseline odometry travel distance error is smaller in this circular motion dataset, while the \gls{cvo} still underestimates the travel distance. 
    
    \subsubsection{Circular Motion Dataset 2: \texttt{DATA\_10\_02}} \label{subsubsec:circularMotionD2_10_02}
    
    In this dataset, the Coyote~3 rover moves in a circular trajectory with a diameter of approximately $2.5-2.8~\mathrm{m}$. The circular motion is performed in a clockwise direction. This rotation direction can be seen in the X--Y trajectory plot in Figure~\ref{fig:circular_motion_dataset_10_02_trajectory}. The trajectory of the rover does not include any complex maneuvers.
    
    The maximum commanded velocity of the rover is $0.2~\mathrm{m/s}$. The X--Y trajectory comparison and the \texttt{z}-position comparison for \texttt{DATA\_10\_02} are shown in Figures~\ref{fig:circular_motion_dataset_10_02_trajectory} and~\ref{fig:circular_motion_dataset_10_02_z_position}, respectively. The translational and rotational \gls{rpe} plots are shown in Figures~\ref{fig:circular_motion_dataset_10_02_RPE_translational} and~\ref{fig:circular_motion_dataset_10_02_RPE_rotational}, respectively. Table~\ref{tab:odometry_metrics_data-10-02} summarizes some key evaluation metrics for this dataset.
    
    \begin{figure}[!htbp]
        \centering
        \includegraphics[width=0.75\textwidth]{media/results/01_flat_floor_terrain/DATA_10_02/Trajectory_With_Heading_vicon_ground_truth_vs_contact_velocity_odometry_vs_baseline_odometry_DATA_10_02.png}
        \caption{X--Y trajectory comparison for the circular motion dataset \texttt{DATA\_10\_02} on the flat floor terrain. The plot compares the trajectories of the baseline odometry and contact velocity odometry against the ground truth.}
        \label{fig:circular_motion_dataset_10_02_trajectory}
    \end{figure}
    
    \begin{figure}[!htbp]
        \centering
        \includegraphics[width=\textwidth]{media/results/01_flat_floor_terrain/DATA_10_02/z_position_comparison_vicon_ground_truth_vs_contact_velocity_odometry_vs_baseline_odometry_DATA_10_02.png}
        \caption{\texttt{z}-position plot of the circular motion dataset \texttt{DATA\_10\_02} on flat floor terrain. The plot compares the ground truth, baseline odometry, and contact velocity odometry along the \texttt{z} axis.}
        \label{fig:circular_motion_dataset_10_02_z_position}
    \end{figure}
    
    \begin{figure}[!htbp]
        \centering
        \includegraphics[width=\textwidth]{media/results/01_flat_floor_terrain/DATA_10_02/RPE_translational_comparison_contact_velocity_odometry_vs_baseline_odometry_DATA_10_02.png}
        \caption{Translational \gls{rpe} comparison for the circular motion dataset \texttt{DATA\_10\_02} on the flat floor terrain. The plot compares the translational \gls{rpe} of the baseline odometry and contact velocity odometry against the ground truth.}
        \label{fig:circular_motion_dataset_10_02_RPE_translational}
    \end{figure}
    
    \begin{figure}[!htbp]
        \centering
        \includegraphics[width=\textwidth]{media/results/01_flat_floor_terrain/DATA_10_02/RPE_rotational_comparison_contact_velocity_odometry_vs_baseline_odometry_DATA_10_02.png}
        \caption{Rotational \gls{rpe} comparison for the circular motion dataset \texttt{DATA\_10\_02} on the flat floor terrain. The plot compares the rotational \gls{rpe} of the baseline odometry and contact velocity odometry against the ground truth.}
        \label{fig:circular_motion_dataset_10_02_RPE_rotational}
    \end{figure}      
    
    \begin{table}[ht]
    \centering
    \caption{Odometry evaluation metrics for DATA\_10\_02.}
    \label{tab:odometry_metrics_data-10-02}
    \begin{tabular}{lll}
    \hline
    Metric & Baseline odometry & Contact velocity odometry \\
    \hline
    Travel distance error & 0.50~m & -0.40~m \\
    Final position error & 0.04~m & 0.16~m \\
    Yaw RMSE & 0.54~$^\circ$ & 0.99~$^\circ$ \\
    Z position RMSE & 0.01~m & 0.05~m \\
    RPE translational drift RMSE & 0.20~m/m & 0.13~m/m \\
    RPE rotational drift RMSE & 5.22~$^\circ$/m & 4.88~$^\circ$/m \\
    \hline
    \end{tabular}
    \end{table}

    
    From the X--Y trajectory plot in Figure~\ref{fig:circular_motion_dataset_10_02_trajectory}, it can be observed that both the baseline odometry and \gls{cvo} trajectories follow the shape of the ground truth trajectory. The translational and rotational \gls{rpe} drift \gls{rmse} are better for the \gls{cvo} compared to the baseline odometry. This can be inferred from Table~\ref{tab:odometry_metrics_data-10-02} and Figures~\ref{fig:circular_motion_dataset_10_02_RPE_translational} and~\ref{fig:circular_motion_dataset_10_02_RPE_rotational}. The \texttt{z}-position plot in Figure~\ref{fig:circular_motion_dataset_09_03_z_position} shows that the \gls{cvo} follows the variations in the ground truth better than the baseline odometry. The \gls{cvo} starts to drift around $20~\mathrm{s}$ after the rover starts moving causing the \texttt{z}-position \gls{rmse} to be higher that the baseline odometry. The travel distance error of the \gls{cvo} is closer to the ground truth travel distance than the baseline odometry. In terms of final position error, the \gls{cvo} is consistent with the previous circular motion dataset but worse than the baseline odometry. 
    
    Overall, the circular motion datasets show that the \gls{cvo} approach achieves slightly lower rotational \gls{rpe} drift than the baseline odometry, although the yaw \gls{rmse} remains lower for the baseline odometry. Even though the \gls{rmse} of yaw is better for the baseline odometry, the \gls{rmse} of rotational \gls{rpe} drift is better for the \gls{cvo} approach. This suggests that the \gls{cvo} approach can provide locally more consistent rotational estimates in circular trajectories, even when its absolute yaw error is higher than the baseline odometry.
       
    \subsection{Complex Trajectories}    
    
    The complex trajectory datasets include the experiments where the rover moves in a complex trajectory that includes turns and maneuvers. The complex trajectory datasets are: \texttt{DATA\_07\_YY}, \texttt{DATA\_08\_YY}, \texttt{DATA\_11\_YY}, \texttt{DATA\_12\_YY}, and \texttt{DATA\_29\_YY}, where \texttt{YY} represents the iteration number. The evaluation results of a representative subset of these datasets is presented in the following subsections.
    
    \subsubsection{Complex Trajectory Dataset 1: \texttt{DATA\_07\_03}} \label{subsubsec:complexTrajectoryD1_07_03}
    
    In this dataset, the rover moves in a figure-\texttt{8} pattern. This includes combination of clockwise and counter-clockwise circular motion. The maximum commanded velocity of the rover is $0.2~\mathrm{m/s}$. The X--Y trajectory comparison and the \texttt{z}-position comparison for \texttt{DATA\_07\_03} are shown in Figures~\ref{fig:complex_trajectory_dataset_07_03_trajectory} and~\ref{fig:complex_trajectory_dataset_07_03_z_position}, respectively. The translational and rotational \gls{rpe} plots are shown in Figures~\ref{fig:complex_trajectory_dataset_07_03_RPE_translational} and~\ref{fig:complex_trajectory_dataset_07_03_RPE_rotational}, respectively. Table~\ref{tab:odometry_metrics_data-07-03} summarizes some key evaluation metrics for this dataset.
    
    \begin{figure}[!htbp]
        \centering
        \includegraphics[width=0.75\textwidth]{media/results/01_flat_floor_terrain/DATA_07_03/Trajectory_With_Heading_vicon_ground_truth_vs_contact_velocity_odometry_vs_baseline_odometry_DATA_07_03.png}
        \caption{X--Y trajectory comparison for the figure-\texttt{8} motion dataset \texttt{DATA\_07\_03} on the flat floor terrain. The plot compares the trajectories of the baseline odometry and contact velocity odometry against the ground truth.}
        \label{fig:complex_trajectory_dataset_07_03_trajectory}
    \end{figure}
    
    \begin{figure}[!htbp]
        \centering
        \includegraphics[width=\textwidth]{media/results/01_flat_floor_terrain/DATA_07_03/z_position_comparison_vicon_ground_truth_vs_contact_velocity_odometry_vs_baseline_odometry_DATA_07_03.png}
        \caption{\texttt{z}-position plot of the figure-\texttt{8} motion dataset \texttt{DATA\_07\_03} on flat floor terrain. The plot compares the ground truth, baseline odometry, and contact velocity odometry along the \texttt{z} axis.}
        \label{fig:complex_trajectory_dataset_07_03_z_position}
    \end{figure}
    
    \begin{figure}[!htbp]
        \centering
        \includegraphics[width=\textwidth]{media/results/01_flat_floor_terrain/DATA_07_03/yaw_error_comparison_vicon_ground_truth_vs_contact_velocity_odometry_vs_baseline_odometry_DATA_07_03.png}
        \caption{Yaw error plot of the figure-\texttt{8} motion dataset \texttt{DATA\_07\_03} on flat floor terrain. The plot compares the ground truth, baseline odometry, and contact velocity odometry along the \texttt{yaw} axis.}
        \label{fig:complex_trajectory_dataset_07_03_yaw_error}
    \end{figure}
    
    
    \begin{figure}[!htbp]
        \centering
        \includegraphics[width=\textwidth]{media/results/01_flat_floor_terrain/DATA_07_03/RPE_translational_comparison_contact_velocity_odometry_vs_baseline_odometry_DATA_07_03.png}
        \caption{Translational \gls{rpe} comparison for the figure-\texttt{8} motion dataset \texttt{DATA\_07\_03} on the flat floor terrain. The plot compares the translational \gls{rpe} of the baseline odometry and contact velocity odometry against the ground truth.}
        \label{fig:complex_trajectory_dataset_07_03_RPE_translational}
    \end{figure}
    
    \begin{figure}[!htbp]
        \centering
        \includegraphics[width=\textwidth]{media/results/01_flat_floor_terrain/DATA_07_03/RPE_rotational_comparison_contact_velocity_odometry_vs_baseline_odometry_DATA_07_03.png}
        \caption{Rotational \gls{rpe} comparison for the figure-\texttt{8} motion dataset \texttt{DATA\_07\_03} on the flat floor terrain. The plot compares the rotational \gls{rpe} of the baseline odometry and contact velocity odometry against the ground truth.}
        \label{fig:complex_trajectory_dataset_07_03_RPE_rotational}
    \end{figure}      
            
        
    \begin{table}[!ht]
    \centering
    \caption{Odometry evaluation metrics for DATA\_07\_03.}
    \label{tab:odometry_metrics_data-07-03}
    \begin{tabular}{lll}
    \hline
    Metric & Baseline odometry & Contact velocity odometry \\
    \hline
    Travel distance error & 0.30~m & -2.70~m \\
    Final position error & 0.24~m & 0.95~m \\
    Yaw RMSE & 0.85~$^\circ$ & 11.51~$^\circ$ \\
    Z position RMSE & 0.05~m & 0.09~m \\
    RPE translational drift RMSE & 0.14~m/m & 0.16~m/m \\
    RPE rotational drift RMSE & 4.12~$^\circ$/m & 4.53~$^\circ$/m \\
    \hline
    \end{tabular}
    \end{table}
    
    From the X--Y trajectory plot shown in Figure~\ref{fig:complex_trajectory_dataset_07_03_trajectory}, it can be observed that the \gls{cvo} trajectory produces the figure-\texttt{8} shape, but it deviates from the ground truth trajectory. The baseline odometry trajectory is closer to the ground truth trajectory in this dataset. The yaw estimation is also poor for the \gls{cvo} compared to the baseline odometry, which can be inferred from the yaw \gls{rmse} in Table~\ref{tab:odometry_metrics_data-07-03} and from the yaw error plot in Figure~\ref{fig:complex_trajectory_dataset_07_03_yaw_error}. 
    
    In this figure-\texttt{8} dataset, translational and rotational \gls{rpe} drift \gls{rmse} of the \gls{cvo} is slightly higher than the baseline odometry. The \texttt{z}-position plot in Figure~\ref{fig:complex_trajectory_dataset_07_03_z_position} shows that the \gls{cvo} follows the oscillation in the ground truth but drifts away from the ground truth over time. The travel distance of the \gls{cvo} is heavily underestimated, which can be inferred from Table~\ref{tab:odometry_metrics_data-07-03}.
    
    \subsubsection{Complex Trajectory Dataset 2: \texttt{DATA\_11\_01}} \label{subsubsec:complexTrajectoryD2_11_01}
    
    In this dataset, the Coyote~3 rover moves in a clockwise square motion. The rover initially moves forward for approximately $4~\mathrm{m}$, then makes an in-place clockwise turn, then moves forward again for approximately $4~\mathrm{m}$. This is repeated until the rover completes a square motion and arrives back at the starting position. This dataset is considered a complex trajectory because it includes both straight-line motion and in-place turns. The maximum commanded velocity of the rover is $0.2~\mathrm{m/s}$ when moving. Since the rover is human operated, the edges of the square are not perfectly straight, and the corners are not perfect right angles. 
    
    The X--Y trajectory comparison and the \texttt{z}-position comparison for \texttt{DATA\_11\_01} are shown in Figures~\ref{fig:complex_trajectory_dataset_11_01_trajectory} and~\ref{fig:complex_trajectory_dataset_11_01_z_position}, respectively. The translational and rotational \gls{rpe} plots are shown in Figures~\ref{fig:complex_trajectory_dataset_11_01_RPE_translational} and~\ref{fig:complex_trajectory_dataset_11_01_RPE_rotational}, respectively. Table~\ref{tab:odometry_metrics_data-11-01} summarizes some key evaluation metrics for this dataset.
    
    \begin{figure}[!htbp]
        \centering
        \includegraphics[width=0.75\textwidth]{media/results/01_flat_floor_terrain/DATA_11_01/Trajectory_With_Heading_vicon_ground_truth_vs_contact_velocity_odometry_vs_baseline_odometry_DATA_11_01.png}
        \caption{X--Y trajectory comparison for the clockwise square motion dataset \texttt{DATA\_11\_01} on the flat floor terrain. The plot compares the trajectories of the baseline odometry and contact velocity odometry against the ground truth.}
        \label{fig:complex_trajectory_dataset_11_01_trajectory}
    \end{figure}
    
    \begin{figure}[!htbp]
        \centering
        \includegraphics[width=\textwidth]{media/results/01_flat_floor_terrain/DATA_11_01/z_position_comparison_vicon_ground_truth_vs_contact_velocity_odometry_vs_baseline_odometry_DATA_11_01.png}
        \caption{\texttt{z}-position plot of the clockwise square motion dataset \texttt{DATA\_11\_01} on flat floor terrain. The plot compares the ground truth, baseline odometry, and contact velocity odometry along the \texttt{z} axis.}
        \label{fig:complex_trajectory_dataset_11_01_z_position}
    \end{figure}
    
    \begin{figure}[!htbp]
        \centering
        \includegraphics[width=\textwidth]{media/results/01_flat_floor_terrain/DATA_11_01/yaw_error_comparison_vicon_ground_truth_vs_contact_velocity_odometry_vs_baseline_odometry_DATA_11_01.png}
        \caption{Yaw error plot of the clockwise square motion dataset \texttt{DATA\_11\_01} on flat floor terrain. The plot compares the ground truth, baseline odometry, and contact velocity odometry along the yaw axis.}
        \label{fig:complex_trajectory_dataset_11_01_yaw_error}
    \end{figure}  
    
    \begin{figure}[!htbp]
        \centering
        \includegraphics[width=\textwidth]{media/results/01_flat_floor_terrain/DATA_11_01/yaw_orientation_comparison_vicon_ground_truth_vs_contact_velocity_odometry_DATA_11_01.png}
        \caption{Yaw orientation plot of the clockwise square motion dataset \texttt{DATA\_11\_01} on flat floor terrain. The plot compares the ground truth, and contact velocity odometry along the yaw axis.}
        \label{fig:complex_trajectory_dataset_11_01_yaw_orientation}
    \end{figure}  
       
    \begin{figure}[!htbp]
        \centering
        \includegraphics[width=\textwidth]{media/results/01_flat_floor_terrain/DATA_11_01/RPE_translational_comparison_contact_velocity_odometry_vs_baseline_odometry_DATA_11_01.png}
        \caption{Translational \gls{rpe} comparison for the clockwise square motion dataset \texttt{DATA\_11\_01} on the flat floor terrain. The plot compares the translational \gls{rpe} of the baseline odometry and contact velocity odometry against the ground truth.}
        \label{fig:complex_trajectory_dataset_11_01_RPE_translational}
    \end{figure}
    
    \begin{figure}[!htbp]
        \centering
        \includegraphics[width=\textwidth]{media/results/01_flat_floor_terrain/DATA_11_01/RPE_rotational_comparison_contact_velocity_odometry_vs_baseline_odometry_DATA_11_01.png}
        \caption{Rotational \gls{rpe} comparison for the clockwise square motion dataset \texttt{DATA\_11\_01} on the flat floor terrain. The plot compares the rotational \gls{rpe} of the baseline odometry and contact velocity odometry against the ground truth.}
        \label{fig:complex_trajectory_dataset_11_01_RPE_rotational}
    \end{figure}      
    
    \begin{table}[ht]
    \centering
    \caption{Odometry evaluation metrics for DATA\_11\_01.}
    \label{tab:odometry_metrics_data-11-01}
    \begin{tabular}{lll}
    \hline
    Metric & Baseline odometry & Contact velocity odometry \\
    \hline
    Travel distance error & -1.48~m & -3.10~m \\
    Final position error & 0.05~m & 0.88~m \\
    Yaw RMSE & 0.43~$^\circ$ & 12.46~$^\circ$ \\
    Z position RMSE & 0.01~m & 0.05~m \\
    RPE translational RMSE & 0.06~m & 0.06~m \\
    RPE translational drift RMSE & 0.22~m/m & 0.25~m/m \\
    RPE rotational RMSE & 0.50~$^\circ$ & 1.23~$^\circ$ \\
    RPE rotational drift RMSE & 1.99~$^\circ$/m & 4.90~$^\circ$/m \\
    \hline
    \end{tabular}
    \end{table}
    
    From the X--Y trajectory plot shown in Figure~\ref{fig:complex_trajectory_dataset_11_01_trajectory}, it can be observed that the \gls{cvo} trajectory produces the square shape, but it deviates from the ground truth trajectory. The yaw estimation of the \gls{cvo} is less accurate compared to the baseline odometry, as indicated by the high yaw RMSE in Table~\ref{tab:odometry_metrics_data-11-01}. This could be the reason for the overall trajectory deviation of the \gls{cvo} from the ground truth. 
    
    % The \gls{cvo} trajectory follows straight line paths better than the point turns, which is evident from translational \gls{rpe} plot in Figure~\ref{fig:complex_trajectory_dataset_11_01_RPE_translational} and the \gls{rpe} translational \gls{rmse} drift in Table~\ref{tab:odometry_metrics_data-11-01}. The point turns can be seen as the change in yaw orientation in the yaw orientation comparison plot in Figure~\ref{fig:complex_trajectory_dataset_11_01_yaw_orientation}. 
    
    The rotational \gls{rpe} shown in Figure~\ref{fig:complex_trajectory_dataset_11_01_RPE_rotational} shows that the \gls{cvo} has a higher rotational error when the rover is performing point turns. This sudden increase in the yaw error during point turns can be observed in the yaw error plot in Figure~\ref{fig:complex_trajectory_dataset_11_01_yaw_error}. The \gls{cvo} estimation also suffers from initial drift in yaw which can be observed in the yaw error plot.
    
    The \texttt{z}-position plot shown in Figure~\ref{fig:complex_trajectory_dataset_11_01_z_position} shows that the error in the \texttt{z}-position estimation decreases when a point turn is performed and increases when the rover starts moving in a straight line. \todo[inline]{why is that?} The \gls{rmse} of the \texttt{z}-position is very high compared to the baseline odometry. The travel distance estimation is heavily underestimated by the \gls{cvo} approach. These results are presented in Table~\ref{tab:odometry_metrics_data-11-01}. 
    
    \subsubsection{Complex Trajectory Dataset 3: \texttt{DATA\_12\_02}} \label{subsubsec:complexTrajectoryD3_12_02}
    
    In this dataset, the Coyote~3 rover moves in a counter-clockwise square motion. Similar to the previous dataset, the rover initially moves forward for approximately $4~\mathrm{m}$, then makes an in-place counter-clockwise turn, then moves forward again for approximately $4~\mathrm{m}$. This is repeated until the rover completes a square motion and arrives back at the starting position. The maximum commanded velocity of the rover is $0.2~\mathrm{m/s}$ when moving. 
    
    The X--Y trajectory comparison and the \texttt{z}-position comparison for \texttt{DATA\_12\_02} are shown in Figures~\ref{fig:complex_trajectory_dataset_12_02_trajectory} and~\ref{fig:complex_trajectory_dataset_12_02_z_position}, respectively. The translational and rotational \gls{rpe} plots are shown in Figures~\ref{fig:complex_trajectory_dataset_12_02_RPE_translational} and~\ref{fig:complex_trajectory_dataset_12_02_RPE_rotational}, respectively. Table~\ref{tab:odometry_metrics_data-12-02} summarizes some key evaluation metrics for this dataset.
    
    \begin{figure}[!htbp]
        \centering
        \includegraphics[width=0.75\textwidth]{media/results/01_flat_floor_terrain/DATA_12_02/Trajectory_With_Heading_vicon_ground_truth_vs_contact_velocity_odometry_vs_baseline_odometry_DATA_12_02.png}
        \caption{X--Y trajectory comparison for the counter-clockwise square motion dataset \texttt{DATA\_12\_02} on the flat floor terrain. The plot compares the trajectories of the baseline odometry and contact velocity odometry against the ground truth.}
        \label{fig:complex_trajectory_dataset_12_02_trajectory}
    \end{figure}
    
    \begin{figure}[!htbp]
        \centering
        \includegraphics[width=\textwidth]{media/results/01_flat_floor_terrain/DATA_12_02/z_position_comparison_vicon_ground_truth_vs_contact_velocity_odometry_vs_baseline_odometry_DATA_12_02.png}
        \caption{\texttt{z}-position plot of the counter-clockwise square motion dataset \texttt{DATA\_12\_02} on flat floor terrain. The plot compares the ground truth, baseline odometry, and contact velocity odometry along the \texttt{z} axis.}
        \label{fig:complex_trajectory_dataset_12_02_z_position}
    \end{figure}
    
    \begin{figure}[!htbp]
        \centering
        \includegraphics[width=\textwidth]{media/results/01_flat_floor_terrain/DATA_12_02/RPE_translational_comparison_contact_velocity_odometry_vs_baseline_odometry_DATA_12_02.png}
        \caption{Translational \gls{rpe} comparison for the counter-clockwise square motion dataset \texttt{DATA\_12\_02} on the flat floor terrain. The plot compares the translational \gls{rpe} of the baseline odometry and contact velocity odometry against the ground truth.}
        \label{fig:complex_trajectory_dataset_12_02_RPE_translational}
    \end{figure}
    
    \begin{figure}[!htbp]
        \centering
        \includegraphics[width=\textwidth]{media/results/01_flat_floor_terrain/DATA_12_02/RPE_rotational_comparison_contact_velocity_odometry_vs_baseline_odometry_DATA_12_02.png}
        \caption{Rotational \gls{rpe} comparison for the counter-clockwise square motion dataset \texttt{DATA\_12\_02} on the flat floor terrain. The plot compares the rotational \gls{rpe} of the baseline odometry and contact velocity odometry against the ground truth.}
        \label{fig:complex_trajectory_dataset_12_02_RPE_rotational}
    \end{figure}          
    
    \begin{table}[ht]
    \centering
    \caption{Odometry evaluation metrics for DATA\_12\_02.}
    \label{tab:odometry_metrics_data-12-02}
    \begin{tabular}{lll}
    \hline
    Metric & Baseline odometry & Contact velocity odometry \\
    \hline
    Travel distance error & -1.62~m & -3.46~m \\
    Final position error & 0.69~m & 0.50~m \\
    Yaw RMSE & 4.98~$^\circ$ & 8.88~$^\circ$ \\
    Z position RMSE & 0.04~m & 0.22~m \\
    RPE translational RMSE & 0.05~m & 0.06~m \\
    RPE translational drift RMSE & 0.21~m/m & 0.23~m/m \\
    RPE rotational RMSE & 0.69~$^\circ$ & 1.11~$^\circ$ \\
    RPE rotational drift RMSE & 2.74~$^\circ$/m & 4.42~$^\circ$/m \\
    \hline
    \end{tabular}
    \end{table}
    
    From the X--Y trajectory plot shown in Figure~\ref{fig:complex_trajectory_dataset_12_02_trajectory}, it can be observed that the \gls{cvo} trajectory produces the square shape, but it deviates from the ground truth trajectory. The \gls{rmse} of yaw is higher for the \gls{cvo} compared to the baseline odometry, which can be inferred from Table~\ref{tab:odometry_metrics_data-12-02}. Similar to the previous dataset, the \gls{cvo} estimation suffers from initial drift in yaw which can be observed in the rotational \gls{rpe} plot shown in Figure~\ref{fig:complex_trajectory_dataset_12_02_RPE_rotational}. The \gls{rmse} of translational \gls{rpe} is slightly higher for the \gls{cvo} compared to the baseline odometry and the \gls{rmse} of rotational \gls{rpe} is significantly higher for the \gls{cvo} compared to the baseline odometry. 
    
    Similar to the previous dataset, the \texttt{z}-position plot shown in Figure~\ref{fig:complex_trajectory_dataset_12_02_z_position} shows that the error in the \texttt{z}-position estimation decreases when a point turn is performed and increases when the rover is starts moving in a straight line. Unlike the previous dataset, the final position error of the \gls{cvo} is better than the baseline odometry but the travel distance error is underestimated and also worse than the baseline odometry. 
    
    \subsubsection{Complex Trajectory Dataset 4: \texttt{DATA\_29\_01}} \label{subsubsec:complexTrajectoryD4_29_01}
    
    In this dataset, the Coyote~3 rover moves in a random exploration type motion. The maximum commanded velocity of the rover is $0.2~\mathrm{m/s}$ when moving. The rover is moved in a random pattern that includes straight-line motion, curved motion, and three-point turning maneuvers. The rover also climbed a small wooden block and a pair of wooden obstacles in this dataset. This trajectory is considered to be a representative of a real-world scenario where the rover is exploring an unknown environment. The rover is human-operated. 
    
    The X--Y trajectory comparison and the \texttt{z}-position comparison for \texttt{DATA\_29\_01} are shown in Figures~\ref{fig:complex_trajectory_dataset_29_01_trajectory} and~\ref{fig:complex_trajectory_dataset_29_01_z_position}, respectively. The translational and rotational \gls{rpe} plots are shown in Figures~\ref{fig:complex_trajectory_dataset_29_01_RPE_translational} and~\ref{fig:complex_trajectory_dataset_29_01_RPE_rotational}, respectively. Table~\ref{tab:odometry_metrics_data-29-01} summarizes some key evaluation metrics for this dataset.
    
    \begin{figure}[!htbp]
        \centering
        \includegraphics[width=0.75\textwidth]{media/results/01_flat_floor_terrain/DATA_29_01/Trajectory_With_Heading_vicon_ground_truth_vs_contact_velocity_odometry_vs_baseline_odometry_DATA_29_01.png}
        \caption{X--Y trajectory comparison for the random exploration motion dataset \texttt{DATA\_29\_01} on the flat floor terrain. The plot compares the trajectories of the baseline odometry and contact velocity odometry against the ground truth.}
        \label{fig:complex_trajectory_dataset_29_01_trajectory}
    \end{figure}
    
    \begin{figure}[!htbp]
        \centering
        \includegraphics[width=\textwidth]{media/results/01_flat_floor_terrain/DATA_29_01/z_position_comparison_vicon_ground_truth_vs_contact_velocity_odometry_vs_baseline_odometry_DATA_29_01.png}
        \caption{\texttt{z}-position plot of the random exploration motion dataset \texttt{DATA\_29\_01} on flat floor terrain. The plot compares the ground truth, baseline odometry, and contact velocity odometry along the \texttt{z} axis.}
        \label{fig:complex_trajectory_dataset_29_01_z_position}
    \end{figure}
    
    \begin{figure}[!htbp]
        \centering
        \includegraphics[width=\textwidth]{media/results/01_flat_floor_terrain/DATA_29_01/RPE_translational_comparison_contact_velocity_odometry_vs_baseline_odometry_DATA_29_01.png}
        \caption{Translational \gls{rpe} comparison for the random exploration motion dataset \texttt{DATA\_29\_01} on the flat floor terrain. The plot compares the translational \gls{rpe} of the baseline odometry and contact velocity odometry against the ground truth.}
        \label{fig:complex_trajectory_dataset_29_01_RPE_translational}
    \end{figure}
    
    \begin{figure}[!htbp]
        \centering
        \includegraphics[width=\textwidth]{media/results/01_flat_floor_terrain/DATA_29_01/RPE_rotational_comparison_contact_velocity_odometry_vs_baseline_odometry_DATA_29_01.png}
        \caption{Rotational \gls{rpe} comparison for the random exploration motion dataset \texttt{DATA\_29\_01} on the flat floor terrain. The plot compares the rotational \gls{rpe} of the baseline odometry and contact velocity odometry against the ground truth.}
        \label{fig:complex_trajectory_dataset_29_01_RPE_rotational}
    \end{figure}      
    
    \begin{table}[ht]
    \centering
    \caption{Odometry evaluation metrics for DATA\_29\_01.}
    \label{tab:odometry_metrics_data-29-01}
    \begin{tabular}{lll}
    \hline
    Metric & Baseline odometry & Contact velocity odometry \\
    \hline
    Travel distance error & 0.82~m & -5.50~m \\
    Final position error & 2.92~m & 1.96~m \\
    Yaw RMSE & 25.78~$^\circ$ & 21.78~$^\circ$ \\
    Z position RMSE & 0.03~m & 0.12~m \\
    RPE translational RMSE & 0.04~m & 0.04~m \\
    RPE translational drift RMSE & 0.14~m/m & 0.16~m/m \\
    RPE rotational RMSE & 1.22~$^\circ$ & 1.28~$^\circ$ \\
    RPE rotational drift RMSE & 4.87~$^\circ$/m & 5.14~$^\circ$/m \\
    \hline
    \end{tabular}
    \end{table}

    From the X--Y trajectory plot shown in Figure~\ref{fig:complex_trajectory_dataset_29_01_trajectory}, it can be observed that the \gls{cvo} produces a trajectory which is similar to the ground truth trajectory. The yaw estimation of the \gls{cvo} is better than the baseline odometry, which can be inferred from the yaw RMSE in Table~\ref{tab:odometry_metrics_data-29-01}. The translational \gls{rpe} drift \gls{rmse} is similar for both methods, while the rotational \gls{rpe} drift \gls{rmse} is slightly higher for the \gls{cvo} approach.
    
    The \texttt{z}-position plot shown in Figure~\ref{fig:complex_trajectory_dataset_29_01_z_position} shows that the error in the \texttt{z}-position estimation heavily drifts over time. The point in trajectory where the rover climbs the pair of wooden obstacles can be observed in the \texttt{z}-position plot around $190~\mathrm{s}$. Both the baseline odometry and \gls{cvo} approaches overestimate the \texttt{z}-position when the rover climbs the wooden obstacles. The travel distance estimation is heavily underestimated by the \gls{cvo} approach compared to a slight overestimation by the baseline odometry. 
   
    Unlike the previous complex trajectory datasets, the yaw \gls{rmse} is better for the \gls{cvo} compared to the baseline odometry. This shows that, in a random exploration type motion, the \gls{cvo} approach can be beneficial, even though the travel distance estimation is heavily underestimated. 
        
    The final position error of the \gls{cvo} is better than the baseline odometry in this dataset. This improvement is due to this specific trajectory and not a representative of the overall performance of the \gls{cvo} approach. The final position error for the \texttt{DATA\_29\_02} dataset is worse for the \gls{cvo} compared to the baseline odometry. The \texttt{DATA\_29\_02} dataset is not presented in this report.
        
    \subsection{Flat Floor Terrain Evaluation Summary} \label{subsec:flatFloorEvalSummary}
    
    The straight-line flat floor experiments show that the baseline odometry performs better than the \gls{cvo} approach in terms of travel distance error and translational \gls{rpe} drift. This shows that, for simple low speed straight-line motion on flat floor terrain, the baseline approach is sufficient to estimate the rover's pose. Even though the \gls{cvo} approach captures the variation pattern of the \texttt{z}-position better in some cases, the overall performance is not better than the baseline odometry.
    
    The circular motion flat floor experiments show that the rotational \gls{rpe} drift is slightly better for the \gls{cvo} approach compared to the baseline odometry. This shows that, in circular motion, the \gls{cvo} approach can provide locally more consistent rotational estimates than the baseline odometry. However, the yaw \gls{rmse} is still lower for the baseline odometry, which indicates that the absolute yaw error is higher for the \gls{cvo} approach.
    
    The complex trajectory experiments show mixed results. For the figure-\texttt{8} and square motion datasets, the baseline odometry performs better than the \gls{cvo} approach in terms of yaw \gls{rmse}, travel distance error, and rotational \gls{rpe} drift. The \gls{cvo} trajectory is able to reproduce the general shape of the motion, but it shows larger yaw errors and severely underestimates the travel distance. In the random exploration dataset, however, the \gls{cvo} approach produces a lower yaw \gls{rmse} than the baseline odometry. This suggests that the proposed method can be useful in irregular trajectories, but its performance is not yet consistent across all flat-floor motion types.
    
    To summarize, the evaluation results on flat floor terrain show that the \gls{cvo} approach can provide better local rotational estimates in circular trajectories, but this improvement is not consistent in complex trajectories. The approach still struggles with yaw estimation and travel distance estimation in several straight-line and structured complex trajectories.
    
    \todo[inline]{To be continued . . .}
    
    \section{Sand Terrain}    

    \subsection{Straight-Line Motion}
    
    \subsection{Circular Motion}

    \subsection{Complex Trajectories}    
    
    
    \section{Crater Terrain}    
    
   
    
\end{document}
